% Local quadratic geometry, illustrated for epsilon>0; 11 pt labels.
\begin{tikzpicture}[x=2.3cm,y=1.1cm,>=Latex,line width=.7pt,
 font=\normalfont\fontsize{11}{13.2}\selectfont]
\fill[gray!10] (-2.05,-1.4) rectangle (2.05,0);
\draw[->] (-2.1,0)--(2.3,0) node[right] {$t$};
\draw[->] (-2,-1.4)--(-2,3.2) node[above] {$h(t)$};
\node[above right] at (1.2,0) {$h=0$: fal};
\node[above right] at (-1.85,2.8) {megengedett oldal: $h>0$};
\draw[blue!65!black,very thick,domain=-1.9:-1,samples=50]
 plot (\x,{\x*\x-1});
\draw[blue!65!black,thick,dashed,domain=-1:1.9,samples=100]
 plot (\x,{\x*\x-1});
\draw[densely dotted] (0,-1.25)--(0,2.7);
\fill (-1,0) circle (1.5pt); \fill (0,-1) circle (1.5pt);
\draw (1,-.06)--(1,.06);
\node[above left,fill=white,inner sep=1pt] at (-1,0) {$t_-$};
\node[above left] at (0,0) {$t_g$};
\node[below right] at (1,0) {$t_+$};
\draw[<->] (.12,0)--(.12,-1) node[midway,right] {$\varepsilon$};
\node[below] at (0,-1.06) {$h(t_g)=-\varepsilon$};
\draw[densely dotted] (-1,-.12)--(-1,-1.85);
\draw[densely dotted] (0,-1.45)--(0,-1.85);
\draw[<->] (-1,-1.75)--(0,-1.75)
 node[midway,below] {$t_g-t_-=\sqrt{2\varepsilon/a}$};
\node[align=center] at (0,-2.8) {szaggatott: ütközés nélküli\\sima továbbvezetés};
\end{tikzpicture}
